\section{Introduction}

An ahead-of-time compiler is typically structured as a frontend, a
collection of optimizations, and a backend.
%
The optimizations in the ``middle-end'' of the compiler are numerous,
time-consuming to develop, hard to get right, and accrete assumptions
about costs that are difficult to excise as hardware platforms evolve.
%
An alternate strategy for implementing some parts of a middle-end is
to use a superoptimizer: a program that looks at the code being
compiled and uses a search procedure, a cost function, and an
equivalence verifier to automatically discover better (or even
optimal) code sequences.
%
The idea dates back at least 37~years~\cite{Fraser79}, and has been
the subject of dozens of papers since then.
%
We created \tool, a synthesizing superoptimizer that automatically
derives novel middle-end optimizations; it was originally designed for
LLVM~\cite{Lattner04} but we have also used it to find new
optimizations for the Microsoft Visual C++ compiler.


Several trends convinced us that it was time to write a new
superoptimizer.
%
There has been increased pressure on compiler
developers due to the adoption of higher-level programming languages and a
proliferation of interesting hardware platforms.
%
SAT and SMT solvers continue to improve; they are already more than
capable of discovering equivalence proofs necessary to verify compiler
optimizations involving tens to hundreds of instructions.
%
Solvers are also a key enabler for program synthesis, which supports
the discovery of new optimizations that are out of reach for naïve
search.
%
Finally, verified compilers appear to be much more difficult to extend
than are traditional compilers.
%
Though we have not yet done so, a natural extension of
superoptimization research would be to use a proof-producing solver to
greatly reduce the effort involved in incorporating new optimizations
into a proved-correct compiler~\cite{Leroy09, Mullen16}.
%
In summary, it appears to be a good time to re-evaluate some aspects
of compiler implementation, including how middle-end optimizers are
constructed.


Our contributions include the design and implementation of \tool{}, a
synthesis-based superoptimizer for a domain-specific intermediate
representation (IR) that resembles a purely functional,
control-flow-free subset of LLVM IR\@.
%
\tool{} implements an extension of Gulwani et al.'s synthesis
algorithm~\cite{Gulwani11}, allowing it to synthesize
LLVM's bitwidth-polymorphic instructions.


\tool{} has two intended use cases.
%
First, its results can be turned into actionable advice for compiler
developers; both LLVM and Microsoft compiler developers have
implemented optimizations suggested by \tool{}.
%
Second, \tool{} can run as an LLVM optimization pass, ensuring that
its code improvements can be exploited by other passes, such as
constant propagation and dead code elimination.
%
When building LLVM itself, \tool{} discovers about 7,900 distinct
optimizations, many of which cannot be performed by LLVM on its own,
and applies them a total of 85,000 times.
%
An \tool-optimized Clang-3.9 binary is almost 3\,MB (4.4\%) smaller
than one built without \tool, though it is also about 2\% slower.
%
(We do not yet know why; in this configuration, the only optimizations
performed by \tool{} were replacing variables by constants. This
should not hurt performance. In fact, versions of \tool{} based on
earlier versions of LLVM did get speedups in this case.)
%
Although the initial compilation of a program using \tool{} is often
$5\times$ to $25\times$ slower than optimized compilation with LLVM,
\tool{}'s discoveries are cached and subsequent compilations have much
lower overhead.
%
For example, the time for \tool{} with a warm cache to compile
LLVM on our test machine is about nine minutes, as opposed to
about eight minutes without \tool{}.
